\documentclass[11pt,a4paper]{appletmanual}
\usepackage[english]{babel}
\manuallang{en}

\appletname{Demand Functions}
\appletsub{Where demand curves come from. Move a price or income, and every
optimal bundle leaves a mark: the trail you leave is the demand function.}
\appletdir{demand-functions}

\begin{document}
\manualtitlepage{Applet manual · Introductory Microeconomics · Universidad de los Andes\\
Translated from a GeoGebra applet.\\
MIT licensed. It opens in a browser; there is nothing to install.}

\manualtoc
\thispagestyle{empty}
\clearpage
\setcounter{page}{1}

% ==========================================================================
\section{What this is}

A demand curve is not a datum. It is the result of solving the consumer's
problem once for every price, and writing the answers down.

In a course that gets said, the finished curve gets drawn, and the class moves
on. This applet does the other thing: the consumer's problem on the left, an
empty pair of axes on the right, and the student moving the price. Every
optimal bundle leaves a mark in the right-hand panel. After a while there is a
curve there, and the student built it point by point.

\figher[1.0]{revealed}{After sweeping $p_x$. The purple points in the left
panel are the price-consumption path; those on the right are the demand curve
just built, with the exact curve revealed underneath.}

\begin{amkey}
The trail and the curve are not two things: the revealed curve passes exactly
through the recorded points. Seeing them coincide is the difference between
knowing the definition of a demand function and understanding it.
\end{amkey}

\subsection{Who it is for}

The demand topic, after the consumer's optimum has been covered. It is the
bridge between the two. It also handles normal, inferior and Giffen goods,
which usually survive as a footnote and here can be manufactured and looked at.

%% Sections shared by every manual, English edition.
%% Pulled in with \input{common-en}. The only thing that differs between
%% applets is \appletfolder, which is also the folder name and the path on the site.

\section{Opening it}

There is nothing to install. The whole application is one \code{index.html}
file: it downloads no libraries, calls no server, and stores nothing outside
your browser.

\subsection{The three ways}

\begin{description}
\item[Double-click the file.] The quickest route. Find
  \code{\appletfolder/index.html} and open it. The browser shows it straight off the
  disk, with an address beginning \code{file://}. It works offline from the
  first moment.
\item[From the web.] If someone has published it, the address ends in
  \code{/\appletfolder/}. Once the page has loaded you can disconnect: it needs the
  network for nothing else.
\item[Served from a folder.] For a class, putting the folder on any static
  server is enough. With Python installed:
  \begin{quote}\code{python3 -m http.server 8000}\end{quote}
  then \code{http://localhost:8000/\appletfolder/} in the browser.
\end{description}

\begin{amnote}
Saving the page with \key{Ctrl}\,+\,\key{S} (or \key{Cmd}\,+\,\key{S} on a
Mac) leaves a copy of your own that still works offline and that can travel on a
memory stick or go out by email.
\end{amnote}

\subsection{Language and theme}

Two pairs of buttons sit at the top right.

\begin{ctrltable}{Control & What it does}
ES / EN & Switches the whole interface between Spanish and English on the spot.
          Nothing is recomputed, so you can switch mid-explanation without
          losing the state you had. \\
Dark / Light & Switches the theme. Dark is the original design; light exists
          for projectors and for printing. \\
\end{ctrltable}

Both choices are remembered in the browser, so next time it opens the way you
left it.

\subsection{Which browser}

Any reasonably recent one: Chrome, Edge, Firefox or Safari, on a computer,
tablet or phone. On a narrow screen the panels stack one above the other
instead of sitting side by side. No account, no permissions, no extensions.


% ==========================================================================
\section{The screen}

\fig[1.0]{interface}{The whole screen. Left: the consumer's problem. Right: the
demand curve being built.}

\begin{description}
\item[Ranges.] The interval considered for each variable. It is both the travel
  of the slider and the stretch the right panel's horizontal axis draws. It is
  the most easily forgotten control and the one that solves the most problems.
\item[Preferences.] Five families, the last one writable.
\item[Parameters.] Those of the chosen family.
\item[Variable being swept.] Which of $p_x$, $p_y$ or $m$ moves. The other two
  are held fixed, and recorded points belong to that one only.
\item[Prices and income.] $p_x$, $p_y$, $m$.
\item[Discover the demand curve.] The buttons that record, sweep and clear.
\item[Layers.] What gets drawn.
\end{description}

\subsection{The two panels}

\begin{description}
\item[Choice · $(x, y)$ plane.] The budget line, the indifference curve through
  the optimal bundle, the bundle itself, and the price-consumption path as it
  forms.
\item[Demand.] The optimal quantity against the variable being moved. Its
  header says which. When the variable is income, the panel becomes the
  \ui{Engel curve}.
\end{description}

The \ui{One panel / Two} switch in the right-hand header separates the demand
for $x$ from that for $y$ into two plots, which is what you need when the two
have very different scales.

\fig[1.0]{engel}{With income as the swept variable, the right panel becomes the
Engel curve.}

\subsection{The dark theme}

\fig[0.95]{dark}{Both panels in the dark theme.}

% ==========================================================================
\section{The five families}

Each brings a closed-form demand, except the custom one, which is solved
numerically.

\begin{ctrltable}{Family & Utility and demand}
Cobb--Douglas & $U = x^{\alpha} y^{1-\alpha}$, with
  $x^\ast = \alpha m/p_x$ and $y^\ast = (1-\alpha) m/p_y$. Each good's demand
  ignores the other good's price: the clean case. \\
CES & $U = (x^{r} + y^{r})^{1/r}$. With $r$ near 1 they are near-perfect
  substitutes; with $r$ strongly negative, near-complements. \\
Linear (substitutes) & $U = a x + y$. Demand is at a corner: all $x$ if
  $a > p_x/p_y$, all $y$ if not. \\
X inferior (Giffen) & $U = \ln(x - \underline{x}) -
  \tfrac{1+b}{b}\ln(\bar y - y)$ with $b = (1-\beta)/\beta$. This is the family
  the Giffen case is built with. \\
Custom & Whatever $U(x, y)$ you type. \\
\end{ctrltable}

\begin{ctrltable}{Parameter & Which family}
$\alpha$ & Cobb--Douglas: $x$'s share of spending. \\
$r$ & CES: the substitution exponent. \\
$a$ & Linear: what $x$ is worth against $y$. \\
$\beta$, $\underline{x}$, $\bar y$ & X inferior: the weight, the subsistence
  minimum of $x$, and the ceiling on $y$. \\
\end{ctrltable}

% ==========================================================================
\section{Discovering the demand curve}

This is the heart of the applet. Four controls:

\begin{ctrltable}{Control & What it does}
Mark while moving & Records a point automatically every time the chosen
  variable moves. The most comfortable setting for a class demonstration. \\
Record point & Saves the current bundle. One point, deliberately. \\
Sweep & Runs the variable from one end of its interval to the other, marking as
  it goes. Press again to stop. \\
Clear & Erases the points for the current variable. \\
\end{ctrltable}

Points belong to the variable that was being moved. Switching from $p_x$ to $m$
does not mix the trails: each keeps its own.

\fig[1.0]{swept}{Part-way through a sweep of $p_x$: the points recorded so far
and the current optimal bundle.}

\subsection{The layers}

\begin{ctrltable}{Layer & What it draws}
Recorded points & The marks left so far. \\
Join the points & Connects them with segments. \\
Reveal the curve & Draws the exact, computed demand function underneath the
  points. This is the check. \\
Price-consumption path & In the $(x, y)$ plane, the trail of optimal
  bundles. \\
Budget line & The current constraint. \\
Indifference curve & The one through the optimal bundle. \\
Demand for x / Demand for y & Which of the two demands the right panel
  draws. \\
Highlight Giffen stretch & Marks in purple the stretches with
  $\partial x^\ast/\partial p_x > 0$. \\
Grid & The background rule. \\
\end{ctrltable}

% ==========================================================================
\section{Ranges and axes}

\begin{amwatch}
The $p_x$, $p_y$ and $m$ sliders are \emph{bounded by the interval} set in the
\ui{Ranges} group. If you try to take $p_y$ below its range minimum it will not
move. Lower the range minimum first. This is the most frequent cause of a
slider that appears not to respond.
\end{amwatch}

The right panel's vertical axis has two modes:

\begin{description}
\item[Fixed vertical axis] (the default). The ceiling is
  $m_{\text{max}}/p_{\text{min}}$, the most that can be bought under the
  current ranges. The advantage is that the axis does not move as you drag, so
  two sweeps can be compared by eye.
\item[Automatic.] The axis fits the curve. You need this when the quantity of
  interest is small next to that ceiling --- the Giffen case, for instance,
  where $x^\ast$ hovers around $0.3$ and the ceiling is $8$.
\end{description}

% ==========================================================================
\section{Normal, inferior, Giffen}

The three readouts on the right settle everything:

\begin{ctrltable}{Readout & What it says}
$\partial x^\ast/\partial p_x$ & The slope of demand. Negative is ordinary;
  positive is Giffen. \\
$\partial x^\ast/\partial m$ & The income effect. Negative means an inferior
  good. \\
Type & \emph{normal} or \emph{inferior}, from the sign of the previous one. \\
\end{ctrltable}

And the verdict at the top names the case:

\begin{description}
\item[x is a normal good here.] Demand falls when its price rises and rises
  with income.
\item[x is an inferior good here.] $\partial x^\ast/\partial m < 0$. Being
  inferior is \emph{necessary} for Giffen, but not sufficient.
\item[x is a Giffen good here.] $\partial x^\ast/\partial p_x > 0$. The income
  effect dominates the substitution effect, which can only happen if $x$ is
  inferior.
\end{description}

\fig[1.0]{giffen}{The Giffen case. Demand rises with the price, and the whole
stretch is highlighted. Look at the left panel: the price-consumption path
climbs almost vertically, because the consumer can barely cut $x$ below their
subsistence minimum.}

\fig[0.85]{giffen-verdict}{The verdict in the Giffen case.}

% ==========================================================================
\section{Classroom activities}

\begin{activity}{Building a demand curve by hand}
\amset Opening values (Cobb--Douglas). \ui{Variable being swept} on $p_x$.
Switch on \ui{Mark while moving}.
\amexp At $p_x = 1.83$: $x^\ast = 1.68$, $y^\ast = 2.11$, spend $5.4$ of $5.4$.
Move $p_x$ slowly from one end to the other. Every bundle leaves a mark on the
right.
\par\smallskip
Once you have fifteen or twenty points, switch on \ui{Reveal the curve}. The
computed curve passes exactly through them. Ask the class what a demand
function is, then: not a datum, but the record of twenty optimisation problems
solved.
\end{activity}

\begin{activity}{Spending does not move}
\amset Cobb--Douglas, sweeping $p_x$.
\amexp The \ui{Spend} readout always says $5.4 / 5.4$, and $y^\ast$ does not
change as $p_x$ moves. Under Cobb--Douglas each good takes a fixed share of
income: $x^\ast = \alpha m / p_x$ does not involve $p_y$, and $y^\ast$ does not
involve $p_x$.
\par\smallskip
Move $\alpha$ and watch the two shares change. It is the property that makes
Cobb--Douglas the exam case, and it is worth seeing before it is assumed.
\end{activity}

\begin{activity}{The Engel curve}
\amset \ui{Variable being swept} on $m$. Sweep.
\amexp The right panel is renamed \ui{Engel curve}. Under Cobb--Douglas it
comes out a straight line through the origin: demand is proportional to income.
\par\smallskip
Switch to \ui{X inferior} and repeat. The Engel curve bends downwards: more
income, less $x$. There is the definition of an inferior good, drawn.
\end{activity}

\begin{activity}{Inferior is not Giffen}
\amset \ui{X inferior} family, opening values.
\amexp $x^\ast = 1.08$, $y^\ast = 3.12$,
$\partial x^\ast/\partial p_x = -0.422$ and
$\partial x^\ast/\partial m = -0.068$. The type is \emph{inferior} and the
verdict says so, but demand still slopes downwards.
\par\smallskip
Ask whether that contradicts anything. It does not: being inferior is necessary
for Giffen, not sufficient. The income effect also has to be strong enough, and
for that the consumer has to be made poor in real terms.
\end{activity}

\begin{activity}{Manufacturing a Giffen good}
\amset \ui{X inferior} family. In \ui{Ranges}, lower the minimum of $p_y$ to
$0.2$. Then set $p_y = 0.35$. Turn off the \ui{Fixed vertical axis}.
\amexp Now $x^\ast = 0.30$, $y^\ast = 10.54$, and
$\partial x^\ast/\partial p_x = +3.1\times 10^{-4}$: \emph{positive}. The
verdict says ``x is a Giffen good here'' and the stretch is highlighted.
\par\smallskip
Sweep $p_x$ and watch the curve rise. What you have done is make the other good
so cheap that $x$ has become the consumer's subsistence good: when $x$ gets
dearer they get poorer, and getting poorer they buy more of the cheap staple.
It is the Irish potato story, and it takes this much setting up to see it.
\end{activity}

\begin{activity}{The two extreme cases}
\amset \ui{Linear (substitutes)} family. Sweep $p_x$.
\amexp Demand is a step: all $x$ until $p_x/p_y$ crosses $a$, then zero at
once. There is no stretch in between.
\par\smallskip
Switch to \ui{CES} and take $r$ down towards $-6$: the curves smooth out until
they resemble complements. The two extremes bracket what a demand curve can do,
and the other families fall between them.
\end{activity}

\begin{activity}{The price-consumption path}
\amset Any family. Switch on \ui{Price-consumption path} and sweep $p_x$.
\amexp In the $(x, y)$ plane the trail of optimal bundles appears. Under
Cobb--Douglas it is horizontal, because $y^\ast$ does not depend on $p_x$.
Under the inferior family it climbs almost vertically.
\par\smallskip
The demand curve in the right panel is that same path, with $p_x$ on the
horizontal axis instead of $x$. Same information, two coordinate systems.
\end{activity}

\begin{activity}{Your own utility function}
\amset \ui{Custom} family. Type \code{x\^{}0.3*y\^{}0.7}, then
\code{min(2x,y)}, then \code{x+4ln(y)}.
\amexp Demand is solved numerically and the trail builds the same way. With
\code{min(2x,y)} the two demands move together; with the quasilinear one, the
demand for $y$ stops depending on income once income passes a threshold.
\par\smallskip
This is the moment to have the class predict the shape before sweeping.
\end{activity}

% ==========================================================================
\section{Expression syntax}

\begin{ctrltable}{Element & What is allowed}
Operators & \code{+ - * / \^{} ( )} \\
Functions & \code{sqrt ln log exp abs min max sin cos} \\
Constants & \code{e}, \code{pi} \\
Implicit multiplication & \code{2xy} is \code{2*x*y} \\
Variables & Only \code{x} and \code{y} \\
\end{ctrltable}

% ==========================================================================
\section{How it is computed}

The four named families carry closed-form demands, derived by hand. The custom
one cannot, so its optimum is found numerically along the budget line: a dense
scan and a golden-section refinement, the same as in the other applets in this
collection. It is robust where a first-order condition finds nothing --- corners
and kinks included.

The derivatives $\partial x^\ast/\partial p_x$ and $\partial x^\ast/\partial m$
in the readouts are central differences on the demand function itself, not
symbolic derivatives: that way they work identically for the closed-form
families and for whatever the user types.

% ==========================================================================
\section{Troubleshooting}

\begin{ctrltable}{Symptom & What is happening}
A price slider will not move & It is hitting its range. Change the minimum or
  maximum under \ui{Ranges}. \\
The curve is flat against the axis & The vertical axis is fixed at
  $m_{\text{max}}/p_{\text{min}}$ and the quantity is small. Turn off
  \ui{Fixed vertical axis}. \\
The points vanish when I change variable & They do not: each variable keeps its
  own trail. Go back and they are there. \\
The right panel is unreadable & The two demands have very different scales.
  Switch to \ui{Two} panels, or turn one of the layers off. \\
I cannot get the Giffen case & It takes all three: the \ui{X inferior} family,
  a low $p_y$ (having lowered its range first), and reading the sign of
  $\partial x^\ast/\partial p_x$, not of $\partial x^\ast/\partial m$. \\
\end{ctrltable}

%% Sections shared by every manual, English edition.
%% Pulled in with \input{common-en}. The only thing that differs between
%% applets is \appletfolder, which is also the folder name and the path on the site.

\section{Publishing it for a class}

Any static host will do: GitHub Pages, Netlify, a university web directory,
even a shared folder served over HTTP. The only thing to upload is the
\code{index.html} file, inside a folder named however you want the address to
read.

The repository ships a GitHub Actions workflow
(\code{.github/workflows/pages.yml}) that publishes every applet at once on each
push to the default branch. To switch it on, in the repository:
\emph{Settings\,$\rightarrow$\,Pages\,$\rightarrow$\,Source: GitHub Actions}.
That is the one step the workflow cannot do for itself, because creating a Pages
site needs administration rights that a \code{GITHUB\_TOKEN} is never granted.

\begin{amnote}
To hand the applet to students without reliable internet, send them the
\code{index.html} file directly. It weighs less than a photograph and opens with
a double-click.
\end{amnote}

\section{Licence and provenance}

The application code is MIT. Use it, change it and pass it on, in class or out
of it, with attribution.

This applet is a standalone rewrite of an original GeoGebra applet. The rewrite
is not a literal translation: where the original had a construction error, a
division by zero reachable with its own sliders, or a dead object inherited from
whatever applet it was copied from, this version fixes it and says so. The
section \emph{Changes from the GeoGebra original} lists those changes one by
one.

Everything is hand-written and dependency-free: the expression parser, the
symbolic differentiation, the level-curve tracing and the drawing. In
particular there is no \code{eval}, so what a student types into a text box
cannot become code, and a strict content security policy does not break the
page.


\end{document}
